% Part I > E. What certificates settle, and what they don't. The problems named
% in section B (and the doubters' external costs), each with what formalization
% does for it. Status and "taken up" columns removed (Kyle, 2026-11-11). Ends
% Part I on the boxed question. 2026-11-11.
% 11% funding figure: Springer Nature survey of 10,480 researchers, via AcadMathSci
% (from the former "Training, access and narrowing" slide).
\section{The Limits of Formalization}

\begin{frame}{What certificates settle, and what they don't}
  \footnotesize\sffamily
  \renewcommand{\arraystretch}{1.3}
  \begin{tabularx}{\linewidth}{@{}p{3.6cm}Y@{}}
    \toprule
    \tabhead{Problem} & \tabhead{What formalization does}\\
    \midrule
    \rowhead{Review overwhelmed} & The kernel checks every step; referees read statements, not proofs. Only where the libraries reach\\
    \rowhead{Undisclosed AI} & The check reads the proof, not the author\\
    \rowhead{Statement fidelity} & A person still has to confirm the formal statement says what was meant\\
    \rowhead{Proof fidelity} & A correct machine proof can still be unreadable, or follow a different argument\\
    \rowhead{Credit and careers} & Registries record what produced a result. Who had the idea is not a formal property\\
    \rowhead{Deskilling and narrowing} & Nothing directly\\
    \rowhead{Funding, access, energy} & Nothing directly. Only 11\% of researchers have institutional funding for their main AI tool\\
    \bottomrule
  \end{tabularx}\par\vspace{0.6cm}
  {\centering\begin{tikzpicture}
    \node[keybox, font=\sffamily\normalsize\bfseries] {What does a principled deployment of Math 2.0 look like?};
  \end{tikzpicture}\par}
  \slidesources{Sections above \textbullet{} Springer Nature survey of 10,480 researchers, via AcadMathSci, Oct 2026}
\end{frame}
